\documentclass{article}

% packages
\usepackage{amsmath, amsthm, thmtools, amsfonts, amssymb, luacode, catchfile, hyperref, ifthen, graphicx, bbm}
\ifcsname c@kobocompile\endcsname
	\usepackage[a5paper, total={1072pt, 1448pt}, margin=10pt, includeheadfoot]{geometry} % set page margins
\else
	\usepackage[a4paper, margin=50pt, includeheadfoot]{geometry}
\fi
\usepackage[shortlabels]{enumitem}
\usepackage[skip=3pt, indent=0pt]{parskip}

% language
\usepackage[bidi=basic, layout=tabular, provide=*]{babel}
\ifcsname c@english\endcsname
	\babelprovide[main, import]{english}
\else
	\babelprovide[main, import, alph=letters.gershayim]{hebrew}
	\setlist[enumerate,2]{label={\alph*.}}
	\babelprovide{rl}
\fi
\babelfont{rm}[Renderer=Harfbuzz]{Libertinus Serif}
\babelfont{sf}{Libertinus Sans}
\babelfont{tt}{Libertinus Mono}

% style
\ifcsname c@stoppagebreaks\endcsname
\else
	\AddToHook{cmd/section/before}{\clearpage}	% Add line break before section
\fi
\linespread{1.3}
\setcounter{secnumdepth}{0}		% Remove default number tags from sections, this won't do well with theorems
\AtBeginDocument{\setlength{\belowdisplayskip}{3pt}}
\AtBeginDocument{\setlength{\abovedisplayskip}{3pt}}
\graphicspath{ {../images/} }

% operators
\DeclareMathOperator\cis{cis}
\DeclareMathOperator\Sp{Span}
\DeclareMathOperator\tr{tr}
\DeclareMathOperator\im{Im}
\DeclareMathOperator\re{Re}
\DeclareMathOperator\diag{diag}
\DeclareMathOperator*\lowlim{\underline{lim}}
\DeclareMathOperator*\uplim{\overline{lim}}
\DeclareMathOperator\rng{rng}
\DeclareMathOperator\Sym{Sym}
\DeclareMathOperator\Arg{Arg}
\DeclareMathOperator\Log{Log}
\DeclareMathOperator\dom{dom}
\DeclareMathOperator\supp{Supp}
\DeclareMathOperator\var{Var}
\DeclareMathOperator\cov{Cov}
\DeclareMathOperator\End{End}
\DeclareMathOperator\id{id}
\DeclareMathOperator\cl{cl}
\DeclareMathOperator\gal{Gal}
\DeclareMathOperator\dist{dist}
\DeclareMathOperator\form{form}
\DeclareMathOperator\term{term}
\DeclareMathOperator\sent{sent}
\DeclareMathOperator\Th{Th}
\DeclareMathOperator\Hom{Hom}
\DeclareMathOperator\Aut{Aut}
\DeclareMathOperator\Char{char}
\DeclareMathOperator\vol{vol}
\DeclareMathOperator\GL{GL}
\DeclareMathOperator\SL{SL}
\DeclareMathOperator\lcm{lcm}

% commands
\newcommand{\cupdot}{\mathbin{\mathaccent\cdot\cup}}
\newcommand{\indicator}{\mathbbm{1}}
\newcommand{\AAA}[0]{\mathbb{A}}
\newcommand{\Aa}[0]{\mathcal{A}}
\newcommand{\BB}[0]{\mathbb{B}}
\newcommand{\Bb}[0]{\mathcal{B}}
\newcommand{\CC}[0]{\mathbb{C}}
\newcommand{\Cc}[0]{\mathcal{C}}
\newcommand{\DD}[0]{\mathbb{D}}
\newcommand{\Dd}[0]{\mathcal{D}}
\newcommand{\EE}[0]{\mathbb{E}}
\newcommand{\FF}[0]{\mathbb{F}}
\newcommand{\Ff}[0]{\mathcal{F}}
\newcommand{\Ii}[0]{\mathcal{I}}
\newcommand{\Gg}[0]{\mathcal{G}}
\newcommand{\HH}[0]{\mathbb{H}}
\newcommand{\Hh}[0]{\mathcal{H}}
\newcommand{\Kk}[0]{\mathcal{K}}
\newcommand{\Ll}[0]{\mathcal{L}}
\newcommand{\Mm}[0]{\mathcal{M}}
\newcommand{\NN}[0]{\mathbb{N}}
\newcommand{\Nn}[0]{\mathcal{N}}
\newcommand{\PP}[0]{\mathbb{P}}
\newcommand{\Pp}[0]{\mathcal{P}}
\newcommand{\QQ}[0]{\mathbb{Q}}
\newcommand{\RR}[0]{\mathbb{R}}
\newcommand{\Rr}[0]{\mathcal{R}}
\newcommand{\Ss}[0]{\mathcal{S}}
\newcommand{\TT}[0]{\mathbb{T}}
\newcommand{\Tt}[0]{\mathcal{T}}
\newcommand{\UU}[0]{\mathbb{U}}
\newcommand{\Uu}[0]{\mathcal{U}}
\newcommand{\Vv}[0]{\mathcal{V}}
\newcommand{\Ww}[0]{\mathcal{W}}
\newcommand{\Xx}[0]{\mathcal{X}}
\newcommand{\Yy}[0]{\mathcal{Y}}
\newcommand{\ZZ}[0]{\mathbb{Z}}
\newcommand{\Zz}[0]{\mathcal{Z}}
\newcommand{\acts}[0]{\circlearrowright}
\newcommand{\explain}[2] {
	\begin{flalign*}
		 && \text{#2} && \text{#1}
	\end{flalign*}
}
\newcommand{\maketitleprint}[1][green]{
	\begin{center}
		\includegraphics[width=6cm]{dragon_#1}
	\end{center}
}

% theorem commands
\newtheoremstyle{c_remark}
	{}	% Space above
	{}	% Space below
	{}% Body font
	{}	% Indent amount
	{\bfseries}	% Theorem head font
	{}	% Punctuation after theorem head
	{.5em}	% Space after theorem head
	{\thmname{#1}\thmnumber{ #2}\thmnote{ \normalfont{\text{(#3)}}}}	% head content
\newtheoremstyle{c_definition}
	{3pt}	% Space above
	{3pt}	% Space below
	{}% Body font
	{}	% Indent amount
	{\bfseries}	% Theorem head font
	{}	% Punctuation after theorem head
	{.5em}	% Space after theorem head
	{\thmname{#1}\thmnumber{ #2}\thmnote{ \normalfont{\text{(#3)}}}}	% head content
\newtheoremstyle{c_plain}
	{3pt}	% Space above
	{3pt}	% Space below
	{\itshape}% Body font
	{}	% Indent amount
	{\bfseries}	% Theorem head font
	{}	% Punctuation after theorem head
	{.5em}	% Space after theorem head
	{\thmname{#1}\thmnumber{ #2}\thmnote{ \text{(#3)}}}	% head content

\ifcsname c@english\endcsname
	\theoremstyle{plain}
	\newtheorem{theorem}{Theorem}[section]
	\newtheorem*{theorem*}{Theorem}
	\newtheorem{lemma}[theorem]{Lemma}
	\newtheorem{proposition}[theorem]{Proposition}
	\newtheorem*{proposition*}{Proposition}

	\theoremstyle{definition}
	\newtheorem{definition}[theorem]{Definition}
	\newtheorem*{definition*}{Definition}
	\newtheorem{example}{Example}[section]
	\newtheorem{exercise}{Exercise}[section]

	\theoremstyle{remark}
	\newtheorem*{remark}{Remark}
	\newtheorem*{solution}{Solution}
	\newtheorem{corollary}[theorem]{Corollary}
	\newtheorem{notation}[theorem]{Notation}
\else
	\theoremstyle{c_plain}
	\newtheorem{theorem}{משפט}[section]
	\newtheorem*{theorem*}{משפט}
	\newtheorem{lemma}[theorem]{למה}
	\newtheorem{proposition}[theorem]{טענה}
	\newtheorem*{proposition*}{טענה}

	\theoremstyle{c_definition}
	\newtheorem{definition}[theorem]{הגדרה}
	\newtheorem*{definition*}{הגדרה}
	\newtheorem{example}{דוגמה}[section]
	\newtheorem{exercise}{תרגיל}[section]

	\theoremstyle{c_remark}
	\newtheorem*{remark}{הערה}
	\newtheorem*{solution}{פתרון}
	\newtheorem{corollary}[theorem]{מסקנה}
	\newtheorem{notation}[theorem]{סימון}
\fi

% Questions related commands
\newcounter{question}
\setcounter{question}{1}
\newcounter{sub_question}
\setcounter{sub_question}{1}

\ifcsname c@english\endcsname
	\newcommand{\question}[1][0]{
		\ifthenelse{#1 = 0}{}{\setcounter{question}{#1}}
		\section{Question \arabic{question}}
		\addtocounter{question}{1}
		\setcounter{sub_question}{1}
	}

	\newcommand{\subquestion}[1][0]{
		\ifthenelse{#1 = 0}{}{\setcounter{sub_question}{#1}}
		\subsection{Part \alph{sub_question}}
		\addtocounter{sub_question}{1}
	}
\else
	\newcommand{\question}[1][0]{
		\ifthenelse{#1 = 0}{}{\setcounter{question}{#1}}
		\section{שאלה \arabic{question}}
		\addtocounter{question}{1}
		\setcounter{sub_question}{1}
	}

	\newcommand{\subquestion}[1][0]{
		\ifthenelse{#1 = 0}{}{\setcounter{sub_question}{#1}}
		\subsection{סעיף \localecounter{letters.gershayim}{sub_question}}
		\addtocounter{sub_question}{1}
	}
\fi

% import lua and start of document
\directlua{common = require ('../common')}

\GetEnv{AUTHOR}

% headers
\author{\AUTHOR}
\date\today

\title{יסודות בתורת ההצגות --- סיכום}
\setcounter{secnumdepth}{2}
% chktex-file 9
% chktex-file 17

\usepackage{fancyhdr}
\pagestyle{fancy}
\renewcommand{\headrulewidth}{0pt}

\begin{document}
\maketitle
\maketitleprint[red]{}

\tableofcontents

\section{שיעור 1 --- 19.10.2026}
תוכנית הקורס הכללית היא מודולים, תורת Wedderburn, תורת הקרקטרים, אינדוקציה של הצגות, הצגות של $S_n$, דואליות Schur-Weyl (הצגות של $S_n$ ושל $U(n)$).
הספרים הם Alperin-Bell: Groups and representations (פרקים 5 ו־6), Peter Webb: Course in finite group representation theory, הספר Fulton-Harris אך הוא מכיל המון ולבסוף גם Lieback-James.
על תורת ההצגות של $S_n$ נלמד ממאמר של Vershik-Okounkov, וטקסטים קלאסיים של Sagan ו־Fulton.

תורת ההצגות היא שילוב של תורת החבורות והאלגברה הלינארית.
ניזכר בפעולה של חבורה על קבוצה,
\begin{definition}[פעולה של חבורה על קבוצה]
	תהי $G$ חבורה ותהי קבוצה $X$, פעולה של $G$ על $X$ היא הומומורפיזם $\rho : G \to \Sym(X)$, כלומר $\rho(g_1 g_2) = \rho(g_1) \rho(g_2)$ לכל $g_1, g_2 \in G$.
\end{definition}
אם ל־$X$ יש מבנה נוסף אז לעתים נרצה שגם הפעולה של $G$ ''תשמור'' על המבנה הזה, במקרה בו נעסוק יהיה ש־$X$ הוא מרחב־וקטורי.
\begin{definition}[הצגה של חבורה]
	תהי חבורה $G$, הזוג $(V, \rho)$ כאשר $V$ מרחב־וקטורי ו־$\rho : G \to \GL(V)$ הומומורפיזם של חבורות.
	באנגלית Representation או Linear Representation.
\end{definition}
\begin{definition}[מימד ההצגה]
	עבור $(V, \rho)$ הצגה, נגדיר את המימד שלה להיות המימד של $V$.
\end{definition}
במילים אחרות לכל $g \in G$ יש העתקה לינארית הפיכה $T_g : V \to V$ המקיימת,
\[
	T_{g_1 \cdot g_2}
	= T_{g_1} \circ T_{g_2}
\]
\begin{definition}[הצגה נאמנה]
	אם $\rho$ חד־חד ערכית אז ההצגה תיקרא נאמנה (faithful).
\end{definition}
אם $V$ סוף־מימדי אז על־ידי בחירת בסיס ל־$V$ ניתן לחשוב על ההצגה כהומומורפיזם $G \to \GL_{\dim V}(\FF)$.
ולמעשה אם ההצגה נאמנה אז אנו יכולים לייצג אותה כחבורה של מטריצות.
\begin{remark}
	המונח ''הצגה'' של חבורה עשוי להתפרש באחת הדרכים הבאות:
	\begin{enumerate}
		\item כהומומורפיזם מהחבורה לחבורת האוטומורפיזמים של אובייקט מתמטי
		\item הומומורפיזם ל־$\GL(V)$ עבור $V$ מרחב־וקטורי
		\item הומומורפיזם ל־$\GL(V)$ כאשר $V$ מרחב וקטורי מעל $\CC$
	\end{enumerate}
\end{remark}
\begin{remark}
	הצגה היא הפרט פעולה על $V$, ובכיוון ההפוך מעל פעולה $G \to \Sym(X)$ ניתן להגדיר הצגה מעל שדה $\FF$ כלשהו;
	נגדיר את,
	\[
		V = \FF^\times = \left\{ \sum_{x \in X} \alpha_x \mid \alpha_x \in \FF, \alpha_x = 0 \text{ a.s.} \right\}
	\]
	ו־$G$ תפעל על־ידי תמורות על איברי הבסיס $X$. \\
	בפרט אם $|X| = n$ אז $G \to S_n \to \GL_n(\FF)$, כל תמורה נשלחת למטריצת תמורה מתאימה.
\end{remark}
\begin{remark}
	אין כמעט קשר ל''הצגה'' (presentation) של חבורה על־ידי יוצרים ויחסים, כלומר מבנה מהצורה $\langle s_1, \ldots, s_k \mid r_1, \ldots, r_m \rangle$.
\end{remark}
\begin{remark}
	לעתים נאמר ש־$V$ הצגה של $G$ במקום $(V, \rho)$.
\end{remark}
נעבור למספר דוגמות של הצגות.
\begin{example}
	נגדיר $G = C_4 = \ZZ / 4 \ZZ$ ונניח ש־$G = \langle \sigma \rangle$, ונניח כי $V = \RR^2$.
	נוכל להגדיר את הפעולה $G \to \GL_2(\RR)$ על־ידי,
	\[
		\sigma \mapsto
		\begin{pmatrix}
			1 & 0 \\
			0 & 1
		\end{pmatrix}
	\]
	כלומר כהצגת הזהות, או על־ידי,
	\[
		\sigma \mapsto
		\begin{pmatrix}
			1 & 0 \\
			0 & -1
		\end{pmatrix}
	\]
	כלומר שיקוף, או אף על־ידי,
	\[
		\sigma \mapsto
		\begin{pmatrix}
			0 & -1 \\
			1 & 0
		\end{pmatrix}
	\]
	כלומר סיבוב בתשעים מעלות.
\end{example}
\begin{example}
	נגדיר $G = D_n$ חבורת סימטריות של מצולע משוכלל בעל $n$ קודקודים.
	נסמן $\sigma$ סיבוב ב־$\frac{2 \pi}{n}$ ו־$\tau$ שיקוף, לכן $G = \langle \sigma, \tau \rangle$.
	נסמן $V = \RR^2$ ונגדיר את,
	\[
		\rho : G \to \GL(V) \cong \GL_2(\RR)
	\]
	על־ידי,
	\[
		\sigma \mapsto
		\begin{pmatrix}
			\cos \frac{2 \pi}{n} & - \sin \frac{2 \pi}{n} \\
			\sin \frac{2 \pi}{n} & \cos \frac{2 \pi}{n}
		\end{pmatrix},
		\qquad
		\tau \mapsto
		\begin{pmatrix}
			1 & 0 \\
			0 & -1
		\end{pmatrix}
	\]
\end{example}
\begin{example}
	$G = \ZZ = \langle \sigma \rangle$ וכן $V = \RR^2$.
	$\ZZ$ חבורה חופשית ולכן ניתן לשלוח את $\sigma$ לכל מטריצה ב־$\GL_2(\RR)$, דוגמה ספציפית תהיה,
	\[
		\sigma \mapsto
		\begin{pmatrix}
			1 & 1 \\
			0 & 1
		\end{pmatrix}
	\]
	ולכן,
	\[
		\sigma^n \mapsto
		\begin{pmatrix}
			1 & n \\
			0 & 1
		\end{pmatrix}
	\]
\end{example}
\begin{example}
	$G = C_2 = \langle \sigma \rangle$ וכן $V = \FF_2^2$ מעל השדה $\FF = \FF_2$.
	נגדיר את,
	\[
		\sigma \mapsto
		\begin{pmatrix}
			0 & 1 \\
			1 & 0
		\end{pmatrix}
	\]
\end{example}
\begin{example}
	$G = S_n$ וכן $\FF = \CC$ ו־$V = \CC^n$ אז נוכל להגדיר את,
	\[
		\rho : S_n \to \GL_n(\CC)
	\]
	על־ידי $\sigma \mapsto M$ עבור איזושהי מטריצה מתאימה $M$.
\end{example}
\begin{definition}[שקילות הצגות]
	שתי הצגות $(V_1, \rho_1)$ ו־$(V_2, \rho_2)$ נקראות שקולות (equivalent) אם קיימת העתקה לינארית הפיכה $T : V_1 \to V_2$ המתחלפת עם הפעולה, כלומר,
	\[
		\rho_2(g) \circ T
		= T \circ \rho_1(g)
		\iff T^{-1} \circ \rho_2(g) = \rho_1(g)
	\]
	במקרה זה נסמן $(V_1, \rho_1) \cong (V_2, \rho_2)$.
\end{definition}
במצב בו המרחבים סוף־מימדיים אם $\rho_i : G \to \GL_n(\FF)$ עבור $i \in \{1, 2\}$ אז קיימת מטריצה הפיכה $B \in \GL_n(\FF)$ כך שלכל $g \in G$ מתקיים,
\[
	B^{-1} \rho_2(g) B = \rho_1(g)
\]
\begin{example}
	כמה הצגות דו־מימדיות מעל $\CC$ יש ל־$C_2$ עד כדי שקילות? \\
	התשובה היא שעל־ידי שימוש בצורת ז'ורדן האיבר שאיננו היחידה נשלח ביחידות לאחת המטריצות הבאות:
	\[
		\begin{pmatrix}
			1 & 0 \\
			0 & 1
		\end{pmatrix}
		\qquad
		\begin{pmatrix}
			-1 & 0 \\
			0 & 1
		\end{pmatrix}
		\qquad
		\begin{pmatrix}
			-1 & 0 \\
			0 & -1
		\end{pmatrix}
	\]
\end{example}
באופן כללי אחד האתגרים בתורת ההצגות הוא מיון כל ההצגות של חבורה עד כדי שקילות.
\begin{definition}[תת־מרחב אינווריאנטי]
	תהי $(V, \rho)$ הצגה של $G$. תת־מרחב־וקטורי $U \subseteq V$ נקרא אינווריאנטי אם לכל $g \in G$ ולכל $u \in U$ מתקיים $g . u \in U$.
\end{definition}
במטריצות לאחר שינוי בסיס כל $g$ תישלח למטריצת בלוקים כאשר הבלוק שמתחת לווקטורים של $U$ יהיו אפסים.
\begin{example}
	בכל הצגה $(V, \rho)$ תתי־המרחבים $\{ 0 \}, V$ הם אינווריאנטיים.
\end{example}
\begin{example}
	ב־$C_4$ פעולת הביסוס על $\RR^2$ אין נוספים,
	גם בדוגמה של $G_n$ אין נוספים.
	בדוגמה של חבורה חופשית $\Sp\{ {(1, 0)}^t \}$ תת־מרחב אינווריאנטי.
\end{example}
\begin{definition}[הצגה אי־פריקה]
	הצגה נקראת אי־פריקה (irreducible) אם אין לה תת־מרחבים אינווריאנטיים לא טריוויאליים.
\end{definition}
\begin{example}
	$C_4$ על $\RR^2$ על־ידי סיבובים הצגה אי־פריקה.
	גם עבור $D_4$ על $\RR^2$ וכן הצגות חד־מימדיות.
\end{example}
בהינתן שתי הצגות של $G$, $(V_1, \rho_1)$ ו־$(V_2, \rho_2)$ מעל שדה משותף $\FF$,
ניתן לבנות הצגה $(V_1 \oplus V_2, \rho_1 \oplus V_2)$ כאשר,
\[
	(\rho_1 \oplus \rho_2)(v_1, v_2)
	= (\rho_1(g)\ v_1, \rho_2(g)\ v_2)
\]
בשפת המטריצות,
\[
	\rho_1 \oplus \rho_2
	\mapsto
	\begin{pmatrix}
		\rho_1(g) & 0 \\
		0 & \rho_2(g)
	\end{pmatrix}
\]
\begin{definition}[הצגה פריקה לחלוטין]
	הצגה $(V, \rho)$ של $G$ נקראת פריקה לחלוטין (completely reducible) אם היא שקולה לסכום ישר של הצגות אי־פריקות.
\end{definition}
\begin{example}
	$C_4$ וסיבוב וכן $D_n$ וסיבוב הן פריקות לחלוטין. \\
	הצגה של $\ZZ$ שהצגנו איננה פריקה לחלוטין וגם לא אי־פריקה, וכן ההצגה של $C_2$ שהצגנו לא פריקה לחלוטין.
\end{example}

\section{שיעור 2 --- 20.10.2026}
\begin{proposition}
	אם $G$ סופית ו־$(V, \rho)$ הצגה אי־פריקה אז $\dim V \le |G|$.
\end{proposition}
\begin{proof}
	יהי $0 \ne v \in V$ ונגדיר $U = \Sp\{ g . v \mid g \in G \}$, אז $\dim U \le |G|$.
	נראה ש־$U$ אינווריאנטי, $v = 1 . u \in U$ ולכן $u \ne 0$ ועל־כן אם $V$ אי־פריקה אז $V = U$.
	אם $U$ מרחב לא ריק ואינווריאנטי ולכן הוא כלל המרחב.
	יהי $u \in U$ ולכן $u = \sum_{g \in G} \alpha_g (g. v)$ עבור סקלרים $\alpha_g \in \FF$,
	לכל $h \in G$,
	\[
		h . u
		= h . \sum_g \alpha_g (g. u)
		= \sum_g \alpha_g (h . g . v)
		\in U
	\]
	שכן $h . g . v = (h . g) . v \in U$.
\end{proof}
\begin{remark}
	נראה בהמשך שאם הצגה אי־פריקה לחלוטין, אז בכל שני פירוקים להצגות אי־פריקות תופענה אותן הצגות אי־פריקות עם אותו ריבוי (עד כדי שקילות).
\end{remark}
\begin{remark}
	בכל הצגה סוף־מימדית יש יחידות של ''גורמי ההרכב'' (גרסה של משפט ז'ורדן־הלדר להצגות).
\end{remark}
\begin{theorem}[משקה (Maschke)]
	אם $G$ סופית ו־$\Char \FF \nmid |G|$ (או $\Char \FF = 0$) אז כל הצגה סוף־מימדית של $G$ היא פריקה לחלוטין.
\end{theorem}
ראינו דוגמה $\rho : C_2 \to \GL_2(\FF_2)$ שאיננה פריקה, וראינו גם שקיימת הצגה $\rho : \ZZ \to \GL_2(\RR)$ שגם היא איננה פריקה לחלוטין.
\begin{proof}
	מספיק להוכיח תת־מרחב וקטורי אינווריאנטי יש משלים ישר אינווריאנטי,
	כלומר שאם $U \le V$ אינווריאנטי אז קיים $W \le V$ אינווריאנטי כך ש־$U \oplus W = V$.

	אם $\FF$ הוא $\RR$ או $\CC$ ויש מכפלה פנימית על $V$ ו־$U \le V$ תת־מרחב וקטורי אינווריאנטי.
	באופן כללי $U^\perp$ לא בהכרח אינווריאנטי, ואין בהכרח בין מכפלה פנימית והצגה. \\
	אם ההצגה מכבדת את המכפלה הפנימית, כלומר פועלת כאיזומטריה:
	\[
		 \forall g \in G \forall v, u \in V,\ 
		 \langle g . u, g . v \rangle = \langle u, v \rangle
	\]
	אז $U^\perp$ גם אינווריאנטי. נראה כי,
	\[
		w \in U^\perp
		\iff \forall u \in U,\ \langle w, u \rangle = 0
		\iff \forall u \in U, g \in G,\ \langle g . w, u \rangle = \langle w, g^{-1} u \rangle = 0
	\]
	זאת שכן $g^{-1} u \in U$ מאינווריאנטיות. לכן $g . w \in U^\perp$. \\
	באופן כללי אפשר מעל $\RR$ או $\CC$ אפשר לבחור מכפלה פנימית שמכבדת את החבורה (נקרא Weylis Unitary trick).
	ניקח מכפלה פנימית כלשהי $\langle , \rangle$ ונגדיר מכפלה פנימית חדשה על־ידי,
	\[
		\langle u, v \rangle_G
		= \frac{1}{|G|} \sum_{g \in G} \langle g u, g v \rangle
	\]
	זו מכפלה פנימית כסכום של מכפלות פנימיות, ו־$\rho$ פועלת כאיזומטריה על המכפלה הפנימית הזו,
	\[
		\langle h . u, h . v \rangle_G
		= \frac{1}{|G|} \sum_{g \in G} \langle g . h . u, g . h . v \rangle
		= \frac{1}{|G|} \sum_{t = g . h \in G} \langle t . u, t . v \rangle
		= \langle u, v \rangle_G
	\]
	נבחין כי הוכחה זו משתנה כתלות בשדה ($\RR$ או $\CC$),
	וכל שמעל $\RR$ כל הצגה סוף־מימדית של חבורה סופית שקולה להצגה אורתוגונלית.
	עוד נובע שמעל $\CC$ כל הצגה שקולה להצגה אוניטרית, כלומר $\rho : G \to O(n) \subseteq \GL_n(\RR)$.

	תהי $G$ סופית ו־$V$ הצגה סוף מימדית, $U \le V$ אינווריאנטי,
	נבחר משלים ישר כלשהו $W \le V$ כך ש־$W \oplus U = V$.
	תהי $P : V \to U$ ההטלה, כלומר $P(u) = u$ לכל $u \in U$, ו־$P(w) = 0$ לכל $w \in W$.
	נגדיר $\bar{P} : V \to V$ על־ידי,
	\[
		\bar{P}(v)
		= \frac{1}{|G|} \sum_{g \in G} g . P(g^{-1} v)
	\]
	נראה ש־$\bar{P}$ הטלה על $U$ וש־$\ker \bar{P}$ הוא אינווריאנטי. \\
	לכל $u \in U$ נראה ש־$\bar{P}(u) = u$, נבחין כי $g^{-1} . u \in U$ מאינווריאנטיות $U$ ולכן,
	\[
		\bar{P}(u)
		= \frac{1}{|G|} \sum_{g \in G} g . P(g^{-1} . u)
		= \frac{1}{|G|} \sum_{g \in G} g . g^{-1} . u
		= \frac{1}{|G|} \sum_{g \in G} u
		= u
	\]
	נקבל ש־$\im \bar{P} \subseteq U$.
	מכיוון ש־$\bar{P}$ כמו בהגדרה אז $U \oplus \ker \bar{P} = V$, ועתה נראה ש־$\ker \bar{P}$ הוא אינווריאנטי.
	אכן לכל $w \in \ker \bar{P}$ צריך להראות ש־$h . w \in \ker \bar{P}$ לכל $h \in G$,
	\begin{multline*}
		\bar{P}(h . w)
		= \frac{1}{|G|} \sum_{g \in G} g . P(g^{-1} . h . w)
		= \frac{1}{|G|} \sum_{g \in G} h . (h^{-1} . g) . P({(h^{-1} g)}^{-1} . w) = \\
		= \frac{1}{|G|} \sum_{t \in G} h . t P(t^{-1} . w)
		= h . \frac{1}{|G|} \sum_{t \in G} t . P(t^{-1}. w)
		= h . 0
		= 0
	\end{multline*}
	זאת שכן $t^{-1} . w \in \ker P$.
\end{proof}

\listoftheorems[title=הגדרות ומשפטים,ignoreall,show={theorem,definition},swapnumber,onlynamed={proposition}]
\end{document}
